\section{Marco Teórico}
\label{sec:marco_teorico}

Esta sección resume los fundamentos ya formalizados en el reporte de ocho agentes
---al que se remite para las definiciones completas de las operaciones $\lambda$,
$\eta$ y $\beta$--- y desarrolla con detalle los dos elementos teóricos propios de
esta extensión: la taxonomía multimodal en la que se ubica el sistema y el problema de
la pista faltante con sus tres métodos de solución.

\subsection{Sistemas multimodales}
\label{sec:multimodal}

Un sistema multimodal percibe y relaciona información de más de una modalidad
sensorial.
La taxonomía de Baltru\v{s}aitis, Ahuja y Morency \cite{baltrusaitis2019multimodal}
organiza los retos técnicos de estos sistemas en cinco categorías:
\emph{representación} (codificar cada modalidad de forma que preserve su contenido),
\emph{traducción} (mapear contenido de una modalidad a otra),
\emph{alineación} (identificar correspondencias entre elementos de modalidades
distintas),
\emph{fusión} (combinar evidencia de varias modalidades para una decisión) y
\emph{co-aprendizaje} (transferir conocimiento entre modalidades).

EAM-TMS ya era multimodal en tres de estos ejes desde su versión por lotes:
representa texto (fastText 300D cuantizado por magnitud
\cite{bojanowski2017fasttext}) e imagen (latente ResNet18 de 64D cuantizado
\cite{he2016resnet}); traduce entre ambas con la memoria hetero-asociativa de
contenido (label $\leftrightarrow$ latente); y alinea pares texto--imagen en el
registro conjunto de la EHAM \cite{morales2024hetero}.
Lo que esta extensión agrega es la \textbf{entrada sensorial continua} y una tercera
modalidad en el borde: la voz.
Conviene precisar el estatus teórico de esa adición: la voz no constituye un dominio
nuevo de memoria, sino una \emph{traducción de frontera} (voz$\to$texto) realizada por
un modelo externo de reconocimiento de habla \cite{radford2023whisper}, exactamente
como los pixeles se traducen a latente por el encoder convolucional.
La política de diseño del proyecto se mantiene: los módulos de frontera
\emph{traducen}; las memorias \emph{deciden}.
Ninguna decisión de asignación, redirección, rechazo o recuperación proviene de los
traductores.

\subsection{Reconocimiento de voz: Whisper}

Whisper \cite{radford2023whisper} es una familia de modelos encoder--decoder de
reconocimiento de habla entrenados con supervisión débil a gran escala (680{,}000
horas), robustos a acento, ruido y lenguaje espontáneo, con capacidad nativa de
\emph{traducción} a inglés desde otros idiomas.
Esta extensión usa la variante \emph{small} mediante la reimplementación
faster-whisper \cite{systran2023fasterwhisper}, que integra el detector de actividad
de voz Silero \cite{silero2021vad}; el laboratorio, sin embargo, lo desactiva
(\texttt{vad\_filter=False}) y segmenta el habla con un VAD por energía con piso de
ruido adaptativo propio, mejor adaptado al flujo continuo por bloques de 30 ms.
El modo \texttt{translate} resuelve un desajuste práctico del sistema: el vocabulario
de las memorias está en inglés (proviene de ConceptNet \cite{speer2017conceptnet} y
fastText en inglés), pero el hablante puede dictar en español; la traducción en el
borde entrega a la memoria pistas en el idioma que sus relaciones conocen.

\subsection{La EAM y la EHAM en una página}

La Memoria Asociativa Entrópica \cite{pineda2021entropic} almacena funciones sobre una
tabla $\mathbf{R} \in \mathbb{N}^{n \times m}$; en su variante ponderada (W-EAM)
\cite{pineda2022weighted} las celdas acumulan frecuencias.
El \emph{registro} ($\lambda$) suma la relación unitaria del patrón; el
\emph{reconocimiento} ($\eta$) es una implicación material relajada (containment con
tolerancia $\xi$) que además produce pesos por característica
$w_i = \mathbf{R}[i, a_i]$; la \emph{recuperación} ($\beta$) muestrea por columna una
distribución modulada por una gaussiana centrada en la pista
\cite{pineda2023imagery}.
La extensión hetero-asociativa (EHAM)
\cite{morales2024hetero,morales2025missing} almacena pares de dominios en un tensor 4D
$\mathbf{R} \in \mathbb{N}^{n \times p \times (m+1) \times (q+1)}$ y define la
\emph{proyección} de una pista izquierda $\mathbf{a}$ con pesos normalizados $w'_i$
sobre el dominio derecho:
\begin{equation}
    \Pi[j, l] = \begin{cases}
        0 & \text{si } \exists\, i : \mathbf{R}[i, j, a_i, l] = 0 \\[4pt]
        \displaystyle\sum_{i=1}^{n} w'_i\, \mathbf{R}[i, j, a_i, l] & \text{en otro caso}
    \end{cases}
    \label{eq:proyeccion}
\end{equation}
La conjunción de anulación implementa \emph{containment} estricto: solo hay masa donde
todas las características de la pista tienen soporte simultáneo.
Sobre estas primitivas ---tomadas sin modificación de la implementación de referencia
de Pineda \& Morales--- se construye todo lo que sigue.

\subsection{El problema de la pista faltante}
\label{sec:pista_faltante}

La proyección~\eqref{eq:proyeccion} no es un patrón: es una \emph{relación} en el
dominio destino, una distribución de posibilidades por cada característica.
Cuando existe una pista del dominio destino, la recuperación $\beta$ la usa para
concentrar el muestreo.
Pero en el caso de uso central de este sistema ---una frase evoca una imagen--- no
existe pista visual alguna: el plano queda \emph{ampliamente indeterminado}.
Este es el \textbf{problema de la pista faltante} (\emph{missing cue problem}),
contribución central del artículo de Morales \& Pineda \cite{morales2025missing}:
la memoria \emph{sabe} que ahí hay un objeto (el plano tiene soporte), pero lo que
tiene no es una figura sino una superposición que debe decodificarse.

El artículo propone tres métodos para extraer un objeto definido del plano.
Los tres comparten la misma \emph{prueba} (\emph{test}): dado un candidato
$\mathbf{b}$ del dominio destino con pesos $\mathbf{u}$, se le retro-proyecta al
dominio de la pista y se mide su distancia ponderada a la pista original,
\begin{equation}
    d(\mathbf{a}, \mathbf{b}) \;=\;
    \frac{\sum_{i} w_i \, \delta_i\!\left(a_i,\; \Pi^{-1}(\mathbf{b})\right)}
         {\sum_{i} w_i},
    \label{eq:test}
\end{equation}
donde $\Pi^{-1}(\mathbf{b})$ denota la proyección del candidato hacia el dominio
fuente y $\delta_i$ es la desviación cuadrática esperada de la columna $i$ respecto
del valor $a_i$ de la pista.
Un candidato es tanto mejor cuanto más se parece, visto de regreso, a la pista que lo
evocó.
Los métodos difieren en cómo generan candidatos:

\begin{description}
    \item[Random samples (RS).] Se toma \emph{una} muestra del plano (una reducción
    aleatoria columna a columna proporcional a la masa) y se acepta sin prueba.
    Es la línea base: barato y sin garantía de consistencia.

    \item[Sample-and-test (ST).] Se generan $k$ muestras del plano (en este sistema
    $k = 127$, el \texttt{sample\_size} de la implementación de referencia) y se
    elige la de menor distancia~\eqref{eq:test}.
    El muestreo propone; la retro-proyección dispone.

    \item[Sample-and-search (SS).] Igual que ST, seguido de un descenso local: se
    explora la vecindad del mejor candidato (variaciones de $\pm 1$ nivel por
    característica con soporte en el plano) aceptando todo movimiento que reduzca la
    distancia, hasta el óptimo local.
    Es el método de mejor desempeño reportado en el artículo, al mayor costo
    computacional.
\end{description}

Un punto de trazabilidad importante para este proyecto: la operación de recuperación
que el sistema EAM-TMS ha usado desde sus primeras versiones
(\texttt{recall\_with\_sampling\_n\_search} de la implementación de referencia)
\emph{es} sample-and-search.
Es decir, el sistema siempre devolvió objetos definidos por el mejor de los tres
métodos; lo que no hacía era \emph{exponer} los otros dos ni exhibir la masa
indeterminada de la que parten.
Esta extensión los implementa como métodos seleccionables y los muestra lado a lado
(Sección~\ref{sec:arq_pista}).

\subsubsection*{El prototipo emergente}

Junto a los tres métodos, esta extensión define una cuarta lectura con valor
pedagógico: el \textbf{prototipo emergente} del plano,
\begin{equation}
    \hat{b}_j \;=\; \arg\max_{l}\; \Pi[j, l],
    \label{eq:prototipo}
\end{equation}
la elección determinista del valor más masivo de cada columna, sin muestreo ni
prueba.
El prototipo es la imagen más literal de la ``plastilina'': cada característica por
separado elige su valor más frecuente, pero nada garantiza que la combinación de esos
máximos corresponda a \emph{ningún} objeto registrado --- las correlaciones entre
características, que el registro conjunto de la EHAM sí almacena, se pierden al leer
columna a columna.
La Sección~\ref{sec:res_pista} muestra que, medido con la prueba~\eqref{eq:test}, el
prototipo es consistentemente \emph{peor} que cualquiera de los tres métodos del
artículo: la evidencia cuantitativa de que la recuperación asociativa no es lectura de
máximos sino búsqueda de consistencia conjunta.
